\section{Direct and Iterative Solvers in Linear System Solving}
\label{sec:solvers}

We consider the linear system $\mathbf{A}\mathbf{x} = \mathbf{b}$ with
$\mathbf{A} \in \mathbb{R}^{n \times n}$ nonsingular. In ostrich this is the
Schur-complement system solved in every Newton iteration
\citep[eq.~(12)]{kucera2026ostrich}:
\begin{equation}
  \underbrace{\left(\mathbf{J}\tilde{\mathbf{M}}^{-1}\mathbf{J}^\top + \mathbf{C}\right)}_{\mathbf{A}}
  \underbrace{\Delta\boldsymbol{\lambda}}_{\mathbf{x}}
  = \underbrace{\frac{1}{h}\left(\mathbf{J}\tilde{\mathbf{M}}^{-1}\mathbf{r}_{\text{dyn}} - \mathbf{r}_{\text{con}}\right)}_{\mathbf{b}} ,
  \label{eq:ostrich-system}
\end{equation}
where $n$ is the number of active constraint rows, i.e.\ the dimension of
$\boldsymbol{\lambda}$. The same matrix $\mathbf{A}$ appears in the adjoint
solve \citep[eq.~(18)]{kucera2026ostrich}. In the linear-algebra sections we
use the generic names $\mathbf{x}$ and $\mathbf{b}$, and $\mathbf{r} =
\mathbf{b} - \mathbf{A}\mathbf{x}$ denotes the residual of the \emph{linear}
system, not the Newton residual $\mathbf{r}(\mathbf{s})$.

A \emph{direct} solver computes $\mathbf{x}$ with a finite, fixed
sequence of operations, usually by factorizing $\mathbf{A}$ into triangular factors.
Up to rounding errors the result is exact. An \emph{iterative} solver instead
improves an approximation $\mathbf{x}_k$ step by step and stops when the residual
$\mathbf{r}_k = \mathbf{b} - \mathbf{A}\mathbf{x}_k$ is small enough.

\subsection{Direct solvers}
\label{sec:direct-solvers}

\paragraph{Gaussian elimination.}
Gaussian elimination applies row operations that reduce $\mathbf{A}$ to an upper
triangular matrix $\mathbf{U}$. The unknowns are then recovered by back substitution
\citep{wiki:gaussian-elimination}. Recording the row operations gives the LU
factorization $\boldsymbol{\Pi}\mathbf{A} = \mathbf{L}\mathbf{U}$. Here
$\mathbf{L}$ is unit lower triangular, and $\boldsymbol{\Pi}$ is a
permutation matrix from row pivoting, which is needed for numerical
stability in general \citep{golub2013matrix,trefethen1997numerical}. The
method works for any nonsingular square matrix. Factorizing costs about
$\tfrac{2}{3}n^3$ floating-point operations. After that, each right-hand side
needs one forward and one back substitution, about $2n^2$ operations in total
\citep{wiki:gaussian-elimination,golub2013matrix}.

\paragraph{Cholesky factorization.}
If $\mathbf{A}$ is symmetric positive definite (SPD), i.e.\ $\mathbf{A} = \mathbf{A}^\top$ and
$\mathbf{d}^\top \mathbf{A} \mathbf{d} > 0$ for all $\mathbf{d} \neq \mathbf{0}$, it can be written as
\begin{equation}
  \mathbf{A} = \mathbf{L} \mathbf{L}^\top ,
  \label{eq:cholesky}
\end{equation}
where $\mathbf{L}$ is lower triangular with a positive diagonal. This factorization
exists and is unique exactly when $\mathbf{A}$ is SPD \citep{wiki:cholesky,golub2013matrix}.
A solve therefore needs an SPD matrix. The system is then solved by two
triangular solves:
\begin{equation}
  \mathbf{L}\mathbf{y} = \mathbf{b} \quad \text{(forward substitution)}, \qquad
  \mathbf{L}^\top \mathbf{x} = \mathbf{y} \quad \text{(back substitution)} .
\end{equation}
Cholesky exploits symmetry and costs about $\tfrac{1}{3}n^3$ operations,
half of Gaussian elimination. It is numerically stable without pivoting
\citep{wiki:cholesky,trefethen1997numerical}. Like LU, it costs $\mathcal{O}(n^2)$ per
additional right-hand side once $\mathbf{L}$ is known. So one factor can be reused
for several solves with the same matrix, such as a forward solve and an
adjoint solve with the same symmetric $\mathbf{A}$.

Table~\ref{tab:direct-cost} summarizes the costs. Both factorizations need
$\mathcal{O}(n^2)$ memory for a dense $\mathbf{A}$. The cubic factorization cost is why direct
solvers are attractive for small or well-structured systems, while iterative
methods are preferred when $n$ is large.

\begin{table}[h]
  \centering
  \begin{tabular}{@{}lccc@{}}
    \toprule
    Method & Requirement on $\mathbf{A}$ & Factorization & Solve per right-hand side \\
    \midrule
    Gaussian elimination (LU) & nonsingular & $\approx \tfrac{2}{3}n^3$ & $\approx 2n^2$ \\
    Cholesky & SPD & $\approx \tfrac{1}{3}n^3$ & $\approx 2n^2$ \\
    \bottomrule
  \end{tabular}
  \caption{Leading-order floating-point operation counts of dense direct solvers
    \citep{golub2013matrix,trefethen1997numerical}.}
  \label{tab:direct-cost}
\end{table}

\subsection{Iterative solvers}
\label{sec:iterative-solvers}

Krylov subspace methods build the iterate $\mathbf{x}_k$ from the affine space
\begin{equation}
  \mathbf{x}_k \in \mathbf{x}_0 + \mathcal{K}_k(\mathbf{A}, \mathbf{r}_0), \qquad
  \mathcal{K}_k(\mathbf{A}, \mathbf{r}_0) = \operatorname{span}\{\mathbf{r}_0, \mathbf{A}\mathbf{r}_0, \dots, \mathbf{A}^{k-1}\mathbf{r}_0\},
\end{equation}
where $\mathbf{r}_0 = \mathbf{b} - \mathbf{A}\mathbf{x}_0$ \citep{saad2003iterative}. They access $\mathbf{A}$ only through
matrix--vector products, so they need neither $\mathbf{A}$ in explicit form nor
$\mathcal{O}(n^2)$ memory, and they can stop as soon as the solution is accurate
enough. Their cost is (number of iterations) $\times$ (cost of one
iteration), and the number of iterations depends on the spectrum of $\mathbf{A}$. Reducing it is the purpose of preconditioning
(Section~\ref{sec:preconditioning}).

\paragraph{Conjugate gradient (CG).}
For SPD $\mathbf{A}$, CG \citep{hestenes1952methods,wiki:cg} chooses the $\mathbf{x}_k$ that
minimizes the $\mathbf{A}$-norm of the error, $\|\mathbf{x} - \mathbf{x}_k\|_{\mathbf{A}}$, over the
Krylov space, using short recurrences: one matrix--vector product per
iteration and a few vectors of storage. Its convergence is bounded through the
condition number $\kappa = \lambda_{\max}(\mathbf{A})/\lambda_{\min}(\mathbf{A})$
\citep{saad2003iterative,trefethen1997numerical}:
\begin{equation}
  \frac{\|\mathbf{x} - \mathbf{x}_k\|_{\mathbf{A}}}{\|\mathbf{x} - \mathbf{x}_0\|_{\mathbf{A}}}
  \le 2 \left( \frac{\sqrt{\kappa} - 1}{\sqrt{\kappa} + 1} \right)^{k} .
  \label{eq:cg-bound}
\end{equation}

\paragraph{Conjugate residual (CR) and its preconditioned form (PCR).}
The conjugate residual method \citep{stiefel1955relaxationsmethoden,saad2003iterative}
uses the same Krylov space but minimizes the norm of the \emph{residual}
instead of the $\mathbf{A}$-norm of the error. Its search directions are
$\mathbf{A}$-orthogonal in the sense $(\mathbf{A}\mathbf{p}_i)^\top (\mathbf{A}\mathbf{p}_j) = 0$ for $i \ne j$,
instead of CG's $\mathbf{p}_i^\top \mathbf{A} \mathbf{p}_j = 0$. With an SPD preconditioner $\mathbf{P} \approx \mathbf{A}$, PCR runs CR on
$\mathbf{P}^{-1}\mathbf{A}\mathbf{x} = \mathbf{P}^{-1}\mathbf{b}$ in the $\mathbf{P}$-inner product
$\langle \mathbf{d}_1, \mathbf{d}_2 \rangle_{\mathbf{P}} = \mathbf{d}_1^\top \mathbf{P} \mathbf{d}_2$, in which
$\mathbf{P}^{-1}\mathbf{A}$ is self-adjoint. The $k$-th iterate minimizes
\begin{equation}
  \|\mathbf{r}_k\|_{\mathbf{P}^{-1}} = \sqrt{\mathbf{r}_k^\top \mathbf{P}^{-1} \mathbf{r}_k}
  \quad \text{over} \quad
  \mathbf{x}_k \in \mathbf{x}_0 + \mathcal{K}_k(\mathbf{P}^{-1}\mathbf{A},\, \mathbf{P}^{-1}\mathbf{r}_0) .
\end{equation}
Because this residual norm decreases monotonically, CR suits solvers that
stop on a residual tolerance or a fixed iteration budget, such as an
inner solver inside a Newton loop \citep{macklin2019nonsmooth}. One
iteration costs one product with $\mathbf{A}$, one application of $\mathbf{P}^{-1}$, three
inner products and five vector updates. CR keeps the extra recurrence
$\mathbf{v} = \mathbf{A}\mathbf{p}$ so that no second product with $\mathbf{A}$ is needed.

Algorithm~\ref{alg:pcr} is the PCR variant used by default in the ostrich
engine: $\mathbf{P}$ is the Jacobi preconditioner $\operatorname{diag}(\mathbf{A})$
(Section~\ref{sec:preconditioning}), $\mathbf{x}_0 = \mathbf{0}$,
$\mathrm{tol} = \mathrm{atol} = 10^{-5}$ and $k_{\max} = 16$. It differs
from the textbook form in three details. The preconditioned
residual $\mathbf{z} = \mathbf{P}^{-1}\mathbf{r}$ is updated by its own recurrence
(line~\ref{alg:pcr:update}) instead of being recomputed. The step sizes fall
back to $0$ if a denominator is not positive, as a safeguard in finite
precision. The stopping test is on the unpreconditioned residual, and it is
evaluated after each iteration, so at least one iteration is always performed. When many independent systems are solved in a batch, all of them
iterate until the slowest one meets its tolerance or the budget $k_{\max}$
is reached.

\begin{algorithm}[h]
  \caption{Preconditioned conjugate residual (PCR), as used in ostrich}
  \label{alg:pcr}
  \begin{algorithmic}[1]
    \Require SPD $\mathbf{A}$, right-hand side $\mathbf{b}$, SPD preconditioner $\mathbf{P}$, initial guess $\mathbf{x}_0$, tolerances $\mathrm{tol}, \mathrm{atol}$, budget $k_{\max}$
    \State $\mathbf{x} \gets \mathbf{x}_0$, \quad $\mathbf{r} \gets \mathbf{b} - \mathbf{A}\mathbf{x}$
    \State $\mathbf{z} \gets \mathbf{P}^{-1}\mathbf{r}$, \quad $\mathbf{t} \gets \mathbf{A}\mathbf{z}$
    \State $\mathbf{p} \gets \mathbf{z}$, \quad $\mathbf{v} \gets \mathbf{t}$ \Comment{$\mathbf{v} = \mathbf{A}\mathbf{p}$}
    \State $\rho \gets \mathbf{z}^\top \mathbf{t}$, \quad $k \gets 0$
    \Repeat
      \State $\mathbf{y} \gets \mathbf{P}^{-1}\mathbf{v}$
      \State $\gamma \gets \rho / (\mathbf{y}^\top \mathbf{v})$ \textbf{if} $\mathbf{y}^\top \mathbf{v} > 0$ \textbf{else} $0$
      \State $\mathbf{x} \gets \mathbf{x} + \gamma \mathbf{p}$, \quad $\mathbf{r} \gets \mathbf{r} - \gamma \mathbf{v}$
      \State $\mathbf{z} \gets \mathbf{z} - \gamma \mathbf{y}$ \Comment{$= \mathbf{P}^{-1}\mathbf{r}$} \label{alg:pcr:update}
      \State $\mathbf{t} \gets \mathbf{A}\mathbf{z}$
      \State $\rho_{\text{new}} \gets \mathbf{z}^\top \mathbf{t}$
      \State $\beta \gets \rho_{\text{new}} / \rho$ \textbf{if} $\rho > 0$ \textbf{else} $0$
      \State $\mathbf{p} \gets \mathbf{z} + \beta \mathbf{p}$, \quad $\mathbf{v} \gets \mathbf{t} + \beta \mathbf{v}$
      \State $\rho \gets \rho_{\text{new}}$, \quad $k \gets k + 1$
    \Until{$\|\mathbf{r}\|_2 \le \max(\mathrm{atol},\, \mathrm{tol}\,\|\mathbf{b}\|_2)$ \textbf{or} $k = k_{\max}$}
    \State \Return $\mathbf{x}$
  \end{algorithmic}
\end{algorithm}
